% =====================================================================
%  Capitolo 6 — Discussione e conclusioni
%  Incluso da main.tex con % =====================================================================
%  Capitolo 6 — Discussione e conclusioni
%  Incluso da main.tex con % =====================================================================
%  Capitolo 6 — Discussione e conclusioni
%  Incluso da main.tex con % =====================================================================
%  Capitolo 6 — Discussione e conclusioni
%  Incluso da main.tex con \input{capitolo6_discussione_conclusioni}
%
%  Non introduce misure nuove: tutti i numeri citati provengono dai
%  Capitoli 3, 4 e 5 e sono riportati con il riferimento alla sezione
%  in cui sono stati ottenuti.
% =====================================================================

\chapter{Discussione e conclusioni}
\label{cap:discussione}

I due capitoli sperimentali rispondono ciascuno a metà della domanda posta nel
§\ref{sec:intro-domanda}, e lo fanno su corpora diversi. Il
Capitolo~\ref{cap:risultati1} stabilisce dove un modello produca linguaggio; il
Capitolo~\ref{cap:risultati2} stabilisce che cosa manchi al testo generato là
dove esso è ben formato. Sono due risultati indipendenti, uno positivo e uno negativo, e restano tali:
i corpora sono diversi e le condizioni di generazione del secondo non sono
note, per cui nessuna delle due misure può essere usata come conferma
dell'altra.

Il capitolo non rimisura nulla: il §\ref{sec:risposta} enuncia i due risultati
e il §\ref{sec:conclusioni} raccoglie le conclusioni.


% ---------------------------------------------------------------------
\section{I due risultati}
\label{sec:risposta}
\label{sec:convergenza-cosa}
\label{sec:persistenza}

La domanda del §\ref{sec:intro-domanda} è: quali proprietà statistiche del
linguaggio un modello riproduce effettivamente, e a quale livello della
gerarchia linguistica smette di riprodurle? I due capitoli sperimentali ne
danno due metà, misurate su corpora diversi e con protocolli diversi.

\paragraph{Esito dello sweep di temperatura.}
Un modello riproduce bene le regolarità che dipendono dalla distribuzione delle
frequenze e dalla struttura locale della successione: la legge di Zipf, la
crescita del vocabolario, la comprimibilità, le correlazioni misurate sui
caratteri. Lo fa però soltanto dentro una finestra stretta. Otto criteri tratti
da quattro lavori indipendenti, costruiti per scopi diversi e mai progettati per
concordare, danno una mediana identica $T = 1.0$ per tutti e cinque i modelli,
senza dipendenza rilevabile dalla dimensione (§\ref{sec:tstar}); la transizione
si consuma quasi interamente fra $T = 0.9$ e $T = 1.0$, dove l'entropia passa da
$0.332$ a $5.77$ bit per carattere e il parametro di periodicità cade di un
fattore cinquanta. Fuori da quella finestra il testo degenera nei due modi
opposti descritti nel §\ref{sec:fallimenti}, e i documenti che sostengono
un'analisi semantica sono due su \num{242}. Il capitolo lascia però aperta la
domanda che lo ha originato, perché gli otto criteri che convergono sono tutti
definiti sui conteggi o sul livello dei caratteri: la temperatura che
individuano è quella a cui il testo generato somiglia di più a quello umano
secondo i livelli su cui la somiglianza è più facile, e nulla in quel capitolo
dice che in quel punto il testo sia organizzato attorno a un contenuto.

\paragraph{Esito della misura di burstiness.}
Il secondo capitolo sperimentale porta la domanda al livello che il primo non
raggiunge, su un corpus in cui la qualità del testo non è una variabile. In
prosa enciclopedica ben formata il testo generato conserva l'apparenza
statistica di quello umano: le parole chiave restano più \emph{bursty} dei
propri controlli a frequenza appaiata anche in Grokipedia, e vi riescono quanto
in Wikipedia, $0.740$ contro $0.692$ (§\ref{sec:tre-corpora}). Cambia
l'intensità. Sulle lettere i quattro corpora sono indistinguibili; il divario si
apre al livello delle parole, dove vale più di un ordine di grandezza rispetto a
quello misurato sui caratteri, e l'entropia degli intervalli lo conferma
eliminando l'effetto della diversa lunghezza delle frasi, con una differenza
appaiata di $-0.105$ bit e significatività inferiore a $10^{-3}$
(§\ref{sec:entropia-intervalli}). Dentro il livello lessicale, però, il divario
non separa le parole di contenuto dai controlli a frequenza appaiata. Il
risultato riguarda un solo sistema generativo, un solo dominio e una sola
lingua, e vale per il testo generato \emph{ben formato}: le condizioni di
campionamento con cui Grokipedia è stata prodotta non sono note, e non è
possibile collocarla sull'asse di temperatura del primo capitolo.


% ---------------------------------------------------------------------
\section{Conclusioni}
\label{sec:conclusioni}

Il lavoro ha misurato le proprietà statistiche del linguaggio generato lungo
due assi, la temperatura di campionamento e la gerarchia dei livelli
linguistici.

\begin{enumerate}
  \item \textbf{Esiste una temperatura privilegiata, e la sua posizione è
    comune.} Otto criteri tratti da quattro lavori indipendenti, mai progettati
    per concordare, danno mediana $T = 1.0$ su tutti e cinque i modelli, senza
    dipendenza rilevabile dalla dimensione. È il risultato empirico più solido
    della tesi.

  \item \textbf{La finestra è stretta e la qualità dell'ottimo non è
    garantita.} La transizione si consuma fra $T = 0.9$ e $T = 1.0$, dove
    l'entropia passa da $0.332$ a $5.77$ bit per carattere. Due modelli su
    cinque portano il MAPE al livello umano e gli altri no: l'esistenza
    dell'ottimo è un fatto generale, la sua profondità una proprietà del
    singolo modello.

  \item \textbf{Il testo generato riproduce le regolarità dei livelli bassi
    della gerarchia e perde terreno salendo.} Sulle lettere i quattro corpora
    sono indistinguibili per tutte e tre le misure impiegate; il divario si apre
    al livello delle parole e vi supera di un ordine di grandezza quello
    misurato sui caratteri. Dentro quel livello, però, non separa le parole di
    contenuto dai controlli a frequenza appaiata.

  \item \textbf{Il deficit è di intensità, non di natura.} Il meccanismo
    descritto in~\cite{altmann2012} è presente anche nella prosa generata e vi
    si decompone nelle stesse proporzioni; le sue parole chiave battono i
    propri controlli a frequenza appaiata quanto quelle di Wikipedia. Manca il
    quanto, non il come.

  \item \textbf{La convergenza dei criteri riguarda i livelli bassi.} Gli otto
    criteri che individuano $T = 1.0$ sono tutti definiti sui conteggi o sui
    caratteri. Se il linguaggio naturale giace sulla superficie critica, come
    sostenuto in~\cite{nakaishi2024}, questo lavoro mostra che vi si arriva
    facendo coincidere le statistiche di quei livelli; che cosa accada ai
    livelli alti in quel punto resta fuori da ciò che lo sweep può misurare,
    perché in quel punto i documenti non degenerati sono due su \num{242}.
\end{enumerate}

I contributi di metodo sono elencati nel §\ref{sec:intro-domanda}. I due
principali sono la stima diretta
dell'esponente di coda e del cut-off, che in~\cite{altmann2012} sono adottati
senza misura e che qui risultano corretti nel valore ma inutilizzabili per
ordinare i corpora, e l'entropia degli intervalli $J$
(§\ref{sec:entropia-intervalli-def}), che non risente della diversa lunghezza
delle frasi ed è circa tre volte più stabile rispetto alla lunghezza di analisi
della misura che sostituisce.

Si torni alla Tabella~\ref{tab:tre-regimi}, da cui il lavoro è partito. Delle
quattro sequenze che vi compaiono una sola è linguaggio, e un unico modello
percorre le altre tre cambiando un solo parametro, sfiorando la seconda in due
documenti su \num{242}. Il testo generato sta fra l'ordine e il disordine, ma
non nel modo in cui vi sta il linguaggio: possiede le regolarità che si
misurano sui conteggi e sui caratteri, e possiede in misura ridotta quella che
nasce dall'essere un discorso su qualcosa. Riprodurre le proprietà statistiche
del linguaggio non equivale a riprodurre il meccanismo che le genera, ed è la
distinzione con cui il §\ref{sec:intro-linguaggio} era partito.

%
%  Non introduce misure nuove: tutti i numeri citati provengono dai
%  Capitoli 3, 4 e 5 e sono riportati con il riferimento alla sezione
%  in cui sono stati ottenuti.
% =====================================================================

\chapter{Discussione e conclusioni}
\label{cap:discussione}

I due capitoli sperimentali rispondono ciascuno a metà della domanda posta nel
§\ref{sec:intro-domanda}, e lo fanno su corpora diversi. Il
Capitolo~\ref{cap:risultati1} stabilisce dove un modello produca linguaggio; il
Capitolo~\ref{cap:risultati2} stabilisce che cosa manchi al testo generato là
dove esso è ben formato. Sono due risultati indipendenti, uno positivo e uno negativo, e restano tali:
i corpora sono diversi e le condizioni di generazione del secondo non sono
note, per cui nessuna delle due misure può essere usata come conferma
dell'altra.

Il capitolo non rimisura nulla: il §\ref{sec:risposta} enuncia i due risultati
e il §\ref{sec:conclusioni} raccoglie le conclusioni.


% ---------------------------------------------------------------------
\section{I due risultati}
\label{sec:risposta}
\label{sec:convergenza-cosa}
\label{sec:persistenza}

La domanda del §\ref{sec:intro-domanda} è: quali proprietà statistiche del
linguaggio un modello riproduce effettivamente, e a quale livello della
gerarchia linguistica smette di riprodurle? I due capitoli sperimentali ne
danno due metà, misurate su corpora diversi e con protocolli diversi.

\paragraph{Esito dello sweep di temperatura.}
Un modello riproduce bene le regolarità che dipendono dalla distribuzione delle
frequenze e dalla struttura locale della successione: la legge di Zipf, la
crescita del vocabolario, la comprimibilità, le correlazioni misurate sui
caratteri. Lo fa però soltanto dentro una finestra stretta. Otto criteri tratti
da quattro lavori indipendenti, costruiti per scopi diversi e mai progettati per
concordare, danno una mediana identica $T = 1.0$ per tutti e cinque i modelli,
senza dipendenza rilevabile dalla dimensione (§\ref{sec:tstar}); la transizione
si consuma quasi interamente fra $T = 0.9$ e $T = 1.0$, dove l'entropia passa da
$0.332$ a $5.77$ bit per carattere e il parametro di periodicità cade di un
fattore cinquanta. Fuori da quella finestra il testo degenera nei due modi
opposti descritti nel §\ref{sec:fallimenti}, e i documenti che sostengono
un'analisi semantica sono due su \num{242}. Il capitolo lascia però aperta la
domanda che lo ha originato, perché gli otto criteri che convergono sono tutti
definiti sui conteggi o sul livello dei caratteri: la temperatura che
individuano è quella a cui il testo generato somiglia di più a quello umano
secondo i livelli su cui la somiglianza è più facile, e nulla in quel capitolo
dice che in quel punto il testo sia organizzato attorno a un contenuto.

\paragraph{Esito della misura di burstiness.}
Il secondo capitolo sperimentale porta la domanda al livello che il primo non
raggiunge, su un corpus in cui la qualità del testo non è una variabile. In
prosa enciclopedica ben formata il testo generato conserva l'apparenza
statistica di quello umano: le parole chiave restano più \emph{bursty} dei
propri controlli a frequenza appaiata anche in Grokipedia, e vi riescono quanto
in Wikipedia, $0.740$ contro $0.692$ (§\ref{sec:tre-corpora}). Cambia
l'intensità. Sulle lettere i quattro corpora sono indistinguibili; il divario si
apre al livello delle parole, dove vale più di un ordine di grandezza rispetto a
quello misurato sui caratteri, e l'entropia degli intervalli lo conferma
eliminando l'effetto della diversa lunghezza delle frasi, con una differenza
appaiata di $-0.105$ bit e significatività inferiore a $10^{-3}$
(§\ref{sec:entropia-intervalli}). Dentro il livello lessicale, però, il divario
non separa le parole di contenuto dai controlli a frequenza appaiata. Il
risultato riguarda un solo sistema generativo, un solo dominio e una sola
lingua, e vale per il testo generato \emph{ben formato}: le condizioni di
campionamento con cui Grokipedia è stata prodotta non sono note, e non è
possibile collocarla sull'asse di temperatura del primo capitolo.


% ---------------------------------------------------------------------
\section{Conclusioni}
\label{sec:conclusioni}

Il lavoro ha misurato le proprietà statistiche del linguaggio generato lungo
due assi, la temperatura di campionamento e la gerarchia dei livelli
linguistici.

\begin{enumerate}
  \item \textbf{Esiste una temperatura privilegiata, e la sua posizione è
    comune.} Otto criteri tratti da quattro lavori indipendenti, mai progettati
    per concordare, danno mediana $T = 1.0$ su tutti e cinque i modelli, senza
    dipendenza rilevabile dalla dimensione. È il risultato empirico più solido
    della tesi.

  \item \textbf{La finestra è stretta e la qualità dell'ottimo non è
    garantita.} La transizione si consuma fra $T = 0.9$ e $T = 1.0$, dove
    l'entropia passa da $0.332$ a $5.77$ bit per carattere. Due modelli su
    cinque portano il MAPE al livello umano e gli altri no: l'esistenza
    dell'ottimo è un fatto generale, la sua profondità una proprietà del
    singolo modello.

  \item \textbf{Il testo generato riproduce le regolarità dei livelli bassi
    della gerarchia e perde terreno salendo.} Sulle lettere i quattro corpora
    sono indistinguibili per tutte e tre le misure impiegate; il divario si apre
    al livello delle parole e vi supera di un ordine di grandezza quello
    misurato sui caratteri. Dentro quel livello, però, non separa le parole di
    contenuto dai controlli a frequenza appaiata.

  \item \textbf{Il deficit è di intensità, non di natura.} Il meccanismo
    descritto in~\cite{altmann2012} è presente anche nella prosa generata e vi
    si decompone nelle stesse proporzioni; le sue parole chiave battono i
    propri controlli a frequenza appaiata quanto quelle di Wikipedia. Manca il
    quanto, non il come.

  \item \textbf{La convergenza dei criteri riguarda i livelli bassi.} Gli otto
    criteri che individuano $T = 1.0$ sono tutti definiti sui conteggi o sui
    caratteri. Se il linguaggio naturale giace sulla superficie critica, come
    sostenuto in~\cite{nakaishi2024}, questo lavoro mostra che vi si arriva
    facendo coincidere le statistiche di quei livelli; che cosa accada ai
    livelli alti in quel punto resta fuori da ciò che lo sweep può misurare,
    perché in quel punto i documenti non degenerati sono due su \num{242}.
\end{enumerate}

I contributi di metodo sono elencati nel §\ref{sec:intro-domanda}. I due
principali sono la stima diretta
dell'esponente di coda e del cut-off, che in~\cite{altmann2012} sono adottati
senza misura e che qui risultano corretti nel valore ma inutilizzabili per
ordinare i corpora, e l'entropia degli intervalli $J$
(§\ref{sec:entropia-intervalli-def}), che non risente della diversa lunghezza
delle frasi ed è circa tre volte più stabile rispetto alla lunghezza di analisi
della misura che sostituisce.

Si torni alla Tabella~\ref{tab:tre-regimi}, da cui il lavoro è partito. Delle
quattro sequenze che vi compaiono una sola è linguaggio, e un unico modello
percorre le altre tre cambiando un solo parametro, sfiorando la seconda in due
documenti su \num{242}. Il testo generato sta fra l'ordine e il disordine, ma
non nel modo in cui vi sta il linguaggio: possiede le regolarità che si
misurano sui conteggi e sui caratteri, e possiede in misura ridotta quella che
nasce dall'essere un discorso su qualcosa. Riprodurre le proprietà statistiche
del linguaggio non equivale a riprodurre il meccanismo che le genera, ed è la
distinzione con cui il §\ref{sec:intro-linguaggio} era partito.

%
%  Non introduce misure nuove: tutti i numeri citati provengono dai
%  Capitoli 3, 4 e 5 e sono riportati con il riferimento alla sezione
%  in cui sono stati ottenuti.
% =====================================================================

\chapter{Discussione e conclusioni}
\label{cap:discussione}

I due capitoli sperimentali rispondono ciascuno a metà della domanda posta nel
§\ref{sec:intro-domanda}, e lo fanno su corpora diversi. Il
Capitolo~\ref{cap:risultati1} stabilisce dove un modello produca linguaggio; il
Capitolo~\ref{cap:risultati2} stabilisce che cosa manchi al testo generato là
dove esso è ben formato. Sono due risultati indipendenti, uno positivo e uno negativo, e restano tali:
i corpora sono diversi e le condizioni di generazione del secondo non sono
note, per cui nessuna delle due misure può essere usata come conferma
dell'altra.

Il capitolo non rimisura nulla: il §\ref{sec:risposta} enuncia i due risultati
e il §\ref{sec:conclusioni} raccoglie le conclusioni.


% ---------------------------------------------------------------------
\section{I due risultati}
\label{sec:risposta}
\label{sec:convergenza-cosa}
\label{sec:persistenza}

La domanda del §\ref{sec:intro-domanda} è: quali proprietà statistiche del
linguaggio un modello riproduce effettivamente, e a quale livello della
gerarchia linguistica smette di riprodurle? I due capitoli sperimentali ne
danno due metà, misurate su corpora diversi e con protocolli diversi.

\paragraph{Esito dello sweep di temperatura.}
Un modello riproduce bene le regolarità che dipendono dalla distribuzione delle
frequenze e dalla struttura locale della successione: la legge di Zipf, la
crescita del vocabolario, la comprimibilità, le correlazioni misurate sui
caratteri. Lo fa però soltanto dentro una finestra stretta. Otto criteri tratti
da quattro lavori indipendenti, costruiti per scopi diversi e mai progettati per
concordare, danno una mediana identica $T = 1.0$ per tutti e cinque i modelli,
senza dipendenza rilevabile dalla dimensione (§\ref{sec:tstar}); la transizione
si consuma quasi interamente fra $T = 0.9$ e $T = 1.0$, dove l'entropia passa da
$0.332$ a $5.77$ bit per carattere e il parametro di periodicità cade di un
fattore cinquanta. Fuori da quella finestra il testo degenera nei due modi
opposti descritti nel §\ref{sec:fallimenti}, e i documenti che sostengono
un'analisi semantica sono due su \num{242}. Il capitolo lascia però aperta la
domanda che lo ha originato, perché gli otto criteri che convergono sono tutti
definiti sui conteggi o sul livello dei caratteri: la temperatura che
individuano è quella a cui il testo generato somiglia di più a quello umano
secondo i livelli su cui la somiglianza è più facile, e nulla in quel capitolo
dice che in quel punto il testo sia organizzato attorno a un contenuto.

\paragraph{Esito della misura di burstiness.}
Il secondo capitolo sperimentale porta la domanda al livello che il primo non
raggiunge, su un corpus in cui la qualità del testo non è una variabile. In
prosa enciclopedica ben formata il testo generato conserva l'apparenza
statistica di quello umano: le parole chiave restano più \emph{bursty} dei
propri controlli a frequenza appaiata anche in Grokipedia, e vi riescono quanto
in Wikipedia, $0.740$ contro $0.692$ (§\ref{sec:tre-corpora}). Cambia
l'intensità. Sulle lettere i quattro corpora sono indistinguibili; il divario si
apre al livello delle parole, dove vale più di un ordine di grandezza rispetto a
quello misurato sui caratteri, e l'entropia degli intervalli lo conferma
eliminando l'effetto della diversa lunghezza delle frasi, con una differenza
appaiata di $-0.105$ bit e significatività inferiore a $10^{-3}$
(§\ref{sec:entropia-intervalli}). Dentro il livello lessicale, però, il divario
non separa le parole di contenuto dai controlli a frequenza appaiata. Il
risultato riguarda un solo sistema generativo, un solo dominio e una sola
lingua, e vale per il testo generato \emph{ben formato}: le condizioni di
campionamento con cui Grokipedia è stata prodotta non sono note, e non è
possibile collocarla sull'asse di temperatura del primo capitolo.


% ---------------------------------------------------------------------
\section{Conclusioni}
\label{sec:conclusioni}

Il lavoro ha misurato le proprietà statistiche del linguaggio generato lungo
due assi, la temperatura di campionamento e la gerarchia dei livelli
linguistici.

\begin{enumerate}
  \item \textbf{Esiste una temperatura privilegiata, e la sua posizione è
    comune.} Otto criteri tratti da quattro lavori indipendenti, mai progettati
    per concordare, danno mediana $T = 1.0$ su tutti e cinque i modelli, senza
    dipendenza rilevabile dalla dimensione. È il risultato empirico più solido
    della tesi.

  \item \textbf{La finestra è stretta e la qualità dell'ottimo non è
    garantita.} La transizione si consuma fra $T = 0.9$ e $T = 1.0$, dove
    l'entropia passa da $0.332$ a $5.77$ bit per carattere. Due modelli su
    cinque portano il MAPE al livello umano e gli altri no: l'esistenza
    dell'ottimo è un fatto generale, la sua profondità una proprietà del
    singolo modello.

  \item \textbf{Il testo generato riproduce le regolarità dei livelli bassi
    della gerarchia e perde terreno salendo.} Sulle lettere i quattro corpora
    sono indistinguibili per tutte e tre le misure impiegate; il divario si apre
    al livello delle parole e vi supera di un ordine di grandezza quello
    misurato sui caratteri. Dentro quel livello, però, non separa le parole di
    contenuto dai controlli a frequenza appaiata.

  \item \textbf{Il deficit è di intensità, non di natura.} Il meccanismo
    descritto in~\cite{altmann2012} è presente anche nella prosa generata e vi
    si decompone nelle stesse proporzioni; le sue parole chiave battono i
    propri controlli a frequenza appaiata quanto quelle di Wikipedia. Manca il
    quanto, non il come.

  \item \textbf{La convergenza dei criteri riguarda i livelli bassi.} Gli otto
    criteri che individuano $T = 1.0$ sono tutti definiti sui conteggi o sui
    caratteri. Se il linguaggio naturale giace sulla superficie critica, come
    sostenuto in~\cite{nakaishi2024}, questo lavoro mostra che vi si arriva
    facendo coincidere le statistiche di quei livelli; che cosa accada ai
    livelli alti in quel punto resta fuori da ciò che lo sweep può misurare,
    perché in quel punto i documenti non degenerati sono due su \num{242}.
\end{enumerate}

I contributi di metodo sono elencati nel §\ref{sec:intro-domanda}. I due
principali sono la stima diretta
dell'esponente di coda e del cut-off, che in~\cite{altmann2012} sono adottati
senza misura e che qui risultano corretti nel valore ma inutilizzabili per
ordinare i corpora, e l'entropia degli intervalli $J$
(§\ref{sec:entropia-intervalli-def}), che non risente della diversa lunghezza
delle frasi ed è circa tre volte più stabile rispetto alla lunghezza di analisi
della misura che sostituisce.

Si torni alla Tabella~\ref{tab:tre-regimi}, da cui il lavoro è partito. Delle
quattro sequenze che vi compaiono una sola è linguaggio, e un unico modello
percorre le altre tre cambiando un solo parametro, sfiorando la seconda in due
documenti su \num{242}. Il testo generato sta fra l'ordine e il disordine, ma
non nel modo in cui vi sta il linguaggio: possiede le regolarità che si
misurano sui conteggi e sui caratteri, e possiede in misura ridotta quella che
nasce dall'essere un discorso su qualcosa. Riprodurre le proprietà statistiche
del linguaggio non equivale a riprodurre il meccanismo che le genera, ed è la
distinzione con cui il §\ref{sec:intro-linguaggio} era partito.

%
%  Non introduce misure nuove: tutti i numeri citati provengono dai
%  Capitoli 3, 4 e 5 e sono riportati con il riferimento alla sezione
%  in cui sono stati ottenuti.
% =====================================================================

\chapter{Discussione e conclusioni}
\label{cap:discussione}

I due capitoli sperimentali rispondono ciascuno a metà della domanda posta nel
§\ref{sec:intro-domanda}, e lo fanno su corpora diversi. Il
Capitolo~\ref{cap:risultati1} stabilisce dove un modello produca linguaggio; il
Capitolo~\ref{cap:risultati2} stabilisce che cosa manchi al testo generato là
dove esso è ben formato. Sono due risultati indipendenti, uno positivo e uno negativo, e restano tali:
i corpora sono diversi e le condizioni di generazione del secondo non sono
note, per cui nessuna delle due misure può essere usata come conferma
dell'altra.

Il capitolo non rimisura nulla: il §\ref{sec:risposta} enuncia i due risultati
e il §\ref{sec:conclusioni} raccoglie le conclusioni.


% ---------------------------------------------------------------------
\section{I due risultati}
\label{sec:risposta}
\label{sec:convergenza-cosa}
\label{sec:persistenza}

La domanda del §\ref{sec:intro-domanda} è: quali proprietà statistiche del
linguaggio un modello riproduce effettivamente, e a quale livello della
gerarchia linguistica smette di riprodurle? I due capitoli sperimentali ne
danno due metà, misurate su corpora diversi e con protocolli diversi.

\paragraph{Esito dello sweep di temperatura.}
Un modello riproduce bene le regolarità che dipendono dalla distribuzione delle
frequenze e dalla struttura locale della successione: la legge di Zipf, la
crescita del vocabolario, la comprimibilità, le correlazioni misurate sui
caratteri. Lo fa però soltanto dentro una finestra stretta. Otto criteri tratti
da quattro lavori indipendenti, costruiti per scopi diversi e mai progettati per
concordare, danno una mediana identica $T = 1.0$ per tutti e cinque i modelli,
senza dipendenza rilevabile dalla dimensione (§\ref{sec:tstar}); la transizione
si consuma quasi interamente fra $T = 0.9$ e $T = 1.0$, dove l'entropia passa da
$0.332$ a $5.77$ bit per carattere e il parametro di periodicità cade di un
fattore cinquanta. Fuori da quella finestra il testo degenera nei due modi
opposti descritti nel §\ref{sec:fallimenti}, e i documenti che sostengono
un'analisi semantica sono due su \num{242}. Il capitolo lascia però aperta la
domanda che lo ha originato, perché gli otto criteri che convergono sono tutti
definiti sui conteggi o sul livello dei caratteri: la temperatura che
individuano è quella a cui il testo generato somiglia di più a quello umano
secondo i livelli su cui la somiglianza è più facile, e nulla in quel capitolo
dice che in quel punto il testo sia organizzato attorno a un contenuto.

\paragraph{Esito della misura di burstiness.}
Il secondo capitolo sperimentale porta la domanda al livello che il primo non
raggiunge, su un corpus in cui la qualità del testo non è una variabile. In
prosa enciclopedica ben formata il testo generato conserva l'apparenza
statistica di quello umano: le parole chiave restano più \emph{bursty} dei
propri controlli a frequenza appaiata anche in Grokipedia, e vi riescono quanto
in Wikipedia, $0.740$ contro $0.692$ (§\ref{sec:tre-corpora}). Cambia
l'intensità. Sulle lettere i quattro corpora sono indistinguibili; il divario si
apre al livello delle parole, dove vale più di un ordine di grandezza rispetto a
quello misurato sui caratteri, e l'entropia degli intervalli lo conferma
eliminando l'effetto della diversa lunghezza delle frasi, con una differenza
appaiata di $-0.105$ bit e significatività inferiore a $10^{-3}$
(§\ref{sec:entropia-intervalli}). Dentro il livello lessicale, però, il divario
non separa le parole di contenuto dai controlli a frequenza appaiata. Il
risultato riguarda un solo sistema generativo, un solo dominio e una sola
lingua, e vale per il testo generato \emph{ben formato}: le condizioni di
campionamento con cui Grokipedia è stata prodotta non sono note, e non è
possibile collocarla sull'asse di temperatura del primo capitolo.


% ---------------------------------------------------------------------
\section{Conclusioni}
\label{sec:conclusioni}

Il lavoro ha misurato le proprietà statistiche del linguaggio generato lungo
due assi, la temperatura di campionamento e la gerarchia dei livelli
linguistici.

\begin{enumerate}
  \item \textbf{Esiste una temperatura privilegiata, e la sua posizione è
    comune.} Otto criteri tratti da quattro lavori indipendenti, mai progettati
    per concordare, danno mediana $T = 1.0$ su tutti e cinque i modelli, senza
    dipendenza rilevabile dalla dimensione. È il risultato empirico più solido
    della tesi.

  \item \textbf{La finestra è stretta e la qualità dell'ottimo non è
    garantita.} La transizione si consuma fra $T = 0.9$ e $T = 1.0$, dove
    l'entropia passa da $0.332$ a $5.77$ bit per carattere. Due modelli su
    cinque portano il MAPE al livello umano e gli altri no: l'esistenza
    dell'ottimo è un fatto generale, la sua profondità una proprietà del
    singolo modello.

  \item \textbf{Il testo generato riproduce le regolarità dei livelli bassi
    della gerarchia e perde terreno salendo.} Sulle lettere i quattro corpora
    sono indistinguibili per tutte e tre le misure impiegate; il divario si apre
    al livello delle parole e vi supera di un ordine di grandezza quello
    misurato sui caratteri. Dentro quel livello, però, non separa le parole di
    contenuto dai controlli a frequenza appaiata.

  \item \textbf{Il deficit è di intensità, non di natura.} Il meccanismo
    descritto in~\cite{altmann2012} è presente anche nella prosa generata e vi
    si decompone nelle stesse proporzioni; le sue parole chiave battono i
    propri controlli a frequenza appaiata quanto quelle di Wikipedia. Manca il
    quanto, non il come.

  \item \textbf{La convergenza dei criteri riguarda i livelli bassi.} Gli otto
    criteri che individuano $T = 1.0$ sono tutti definiti sui conteggi o sui
    caratteri. Se il linguaggio naturale giace sulla superficie critica, come
    sostenuto in~\cite{nakaishi2024}, questo lavoro mostra che vi si arriva
    facendo coincidere le statistiche di quei livelli; che cosa accada ai
    livelli alti in quel punto resta fuori da ciò che lo sweep può misurare,
    perché in quel punto i documenti non degenerati sono due su \num{242}.
\end{enumerate}

I contributi di metodo sono elencati nel §\ref{sec:intro-domanda}. I due
principali sono la stima diretta
dell'esponente di coda e del cut-off, che in~\cite{altmann2012} sono adottati
senza misura e che qui risultano corretti nel valore ma inutilizzabili per
ordinare i corpora, e l'entropia degli intervalli $J$
(§\ref{sec:entropia-intervalli-def}), che non risente della diversa lunghezza
delle frasi ed è circa tre volte più stabile rispetto alla lunghezza di analisi
della misura che sostituisce.

Si torni alla Tabella~\ref{tab:tre-regimi}, da cui il lavoro è partito. Delle
quattro sequenze che vi compaiono una sola è linguaggio, e un unico modello
percorre le altre tre cambiando un solo parametro, sfiorando la seconda in due
documenti su \num{242}. Il testo generato sta fra l'ordine e il disordine, ma
non nel modo in cui vi sta il linguaggio: possiede le regolarità che si
misurano sui conteggi e sui caratteri, e possiede in misura ridotta quella che
nasce dall'essere un discorso su qualcosa. Riprodurre le proprietà statistiche
del linguaggio non equivale a riprodurre il meccanismo che le genera, ed è la
distinzione con cui il §\ref{sec:intro-linguaggio} era partito.
